\section{U007 Precise faults and ordering across asynchronous engines}
\label{sec:u007}

The useful correctness contract for a fused accelerator must say which effects have occurred, which state can be inspected, and which operations can safely be repeated. A matrix operation can finish after a younger scalar instruction, a vector consumer can start before an unrelated transfer finishes, and a remote receiver can consume an output before the sending host observes completion. These are ordinary consequences of overlap. The difficulty arises when a programmer, debugger, or recovery mechanism treats all these events as one undifferentiated notion of completion. The design objective is therefore to preserve necessary causal edges while giving faults a well-defined attribution and recovery boundary. Serializing every engine makes some executions easier to inspect, but discards much of the overlap that U001 and U006 seek to expose.

Consider a distributed tile pipeline. A local matrix stage produces an intermediate, a vector stage transforms it, a communication engine transfers the result, and a remote stage consumes it. The algorithm determines the producer--consumer relation. The ISA or device API determines how that relation becomes an ordering or completion guarantee. The compiler translates the relation into commands and synchronization; the runtime determines submission, participant identity, and failure handling; hardware tracks the resulting state. A missing edge can originate in any of these layers. Conversely, moving an edge from explicit software to dependency-tracking hardware changes its implementation and ownership, without eliminating the edge or the resources needed to enforce it.

\subsection{Four contracts that must not be conflated}

A \emph{precisely attributed fault} identifies the operation that caused a failure, possibly including an instruction address, tile identifier, engine, memory address, and source location. A \emph{precise recoverable state} is stronger: the saved architectural state corresponds to a specified instruction boundary, and execution can resume according to a stated restart rule. The classical precise-interrupt definition is based on a sequential architectural boundary; in-order completion is one implementation, while retained versions and recovery structures permit execution and completion to be less restricted \citep{u007_smith_pleszkun}. For a modern accelerator, the definition must identify its recovery domain and the status of pending effects outside that domain.

A \emph{memory-order contract} specifies permissible observations and synchronization relationships, including address spaces, atomicity, participants, and scopes. It does not necessarily impose one total order over the machine. A \emph{completion and ownership contract} specifies when an operation has reached a named milestone and when a resource may change owners. An operation can have finished reading its source while its destination remains unavailable to a consumer. A receiver can have received a tile while still retaining it for computation. These cases require different events even when one implementation happens to report several of them together.

\begin{table}[tbp]
\centering
\small
\begin{tabularx}{\linewidth}{@{}p{0.21\linewidth}XX@{}}
\toprule
Contract & Question answered & What does not follow automatically \\
\midrule
Fault attribution & Which operation caused this error? & Registers and memory form a restartable snapshot \\
Precise local state & At which specified boundary can this domain resume? & Remote consumers can undo escaped effects \\
Memory ordering & Which observations must precede which others, for which observers? & All transfers are complete or all participants have arrived \\
Completion & Which operation milestone has been reached? & Every consumer has stopped using the payload \\
Ownership release & Who may overwrite or reclaim this storage now? & The containing pool can be deallocated \\
Trace or replay & Which events were recorded or can be reproduced? & Every race or unrecorded schedule has been covered \\
\bottomrule
\end{tabularx}
\caption{Correctness and observability contracts answer different questions. Each guarantee needs an explicit scope.}
\label{tab:u007-contracts}
\end{table}

An asynchronous error adds a further distinction between the \emph{cause} and the \emph{reporting point}. For example, the CUDA driver documents that synchronization-related calls can report errors from earlier asynchronous launches \citep{u007_cuda_context}. The call returning an error is thus a useful observation boundary, but need not be the offending operation. A debugger likewise needs to distinguish the faulting program counter from the later counter at which an engine stopped. Even a correct fault address cannot reconstruct a tensor already overwritten by younger work.

\subsection{A publication and reuse example}

Let device A produce tile $X_e$ for epoch $e$ in source slot $S$, and let device B receive it into slot $D$. Define events $W_e$ for the final producer write, $I_e$ for transfer issue, $S_e$ for source-read completion, $D_e$ for destination arrival, $P_e$ for publication of readiness, $C_e$ for the end of the consumer's reads, and $A_e$ for acknowledgment that the destination slot can be reused. The following is an abstract protocol requirement, not a vendor instruction sequence:
\begin{align}
 W_e &\prec I_e, & S_e &\prec \operatorname{overwrite}(S), \label{eq:u007-source}\\
 D_e &\prec P_e \prec \operatorname{read}_B(D), &
 C_e &\prec A_e \prec \operatorname{overwrite}_{e+1}(D). \label{eq:u007-dest}
\end{align}
The symbol $\prec$ means an established causal relation under the selected API and memory model. It is not a claim that an arbitrary local fence implements that relation. Publication may need a release/acquire pair, a DMA completion event, an ordered data-and-signal primitive, or more than one of these. The choice depends on the path by which payload and signal reach the consumer.

Suppose the payload travels on one channel and the ready flag on another. Source-code issue order alone is insufficient unless the selected interface orders those channels. B can observe the flag while reading old or partially updated data. Adding a local fence with a scope that excludes B also fails to establish the required edge. In another failure, A reuses $S$ after merely enqueuing the transfer, so the transfer observes a mixture of two epochs. In a third, A correctly waits for delivery but overwrites $D$ before B finishes consuming it. All three are synchronization errors, but they require different repairs: publication ordering, source completion, and destination ownership respectively.

The protocol need not place a global barrier after every tile. A finite set of destination slots can carry distinct epochs, and the sender can advance whenever it owns a free slot and the required completion state. This retains overlap while making the ownership proof explicit. It also reveals why an epoch or generation identifier is useful: a delayed readiness signal for $e$ must not satisfy a wait for $e+1$. Whether a finite identifier can wrap safely depends on the maximum lifetime of outstanding work and the reuse protocol; a wide counter by itself is not a proof.

A communication API may permit source release before the receiver can consume the destination. Pallas remote DMA exposes separate send and receive waits \citep{u007_pallas_distributed}. Tenstorrent's documented on-chip NoC path similarly distinguishes source-flush completion from destination-arrival completion, while its low-level Ethernet path requires a higher-level acknowledgment protocol for comparable guarantees \citep{u007_tt_ethernet}. These mechanisms instantiate different subsets of the abstract events. They should not be collapsed into one generic barrier in a performance diagram.

\subsection{What a local precise-state design must establish}

For a proposed superscalar NPU, take a local architectural instruction sequence $i_0,i_1,\ldots$ and a retirement frontier $k$. A precise exception at $i_k$ requires a state consistent with the specified boundary before that instruction, subject to its documented fault or partial-completion semantics. All older architectural effects must be represented as completed effects or irrevocably committed pending effects according to the contract. Younger effects must be invisible to observers covered by the rollback guarantee. This is a local sequential boundary, not a promise of globally sequentially consistent memory.

A sufficient construction for \emph{faults detected before commit} can be stated with four invariants. First, committed register and tile mappings identify the versions produced by the retired prefix. Second, speculative versions cannot overwrite the only recoverable version of any live architectural value. Third, externally observable side effects are either withheld until their issuing operation is irrevocable or covered by an explicitly specified transactional recovery protocol. Fourth, every outstanding operation retains a generation-valid identity and reports either successful completion or a fault before it becomes eligible to commit. An implementation may enforce these invariants with renaming, checkpoints, result buffers, ordered publication, or a mixture; a ROB is only one part of the construction.

The corresponding argument is inductive. Initially the committed mappings and effects represent the empty prefix. A successful retirement updates them only after the instruction's required results and fault checks are complete, preserving the next prefix. If the oldest faulting instruction reaches the frontier, younger speculative versions are discarded and the previous committed mappings restored. The side-effect invariant prevents discarded instructions from leaving an observable effect inside the promised recovery domain. The identity invariant prevents a delayed response from an old execution from updating a recycled destination. Under these assumptions the recovered state represents the last retired prefix. This argument is conditional on the invariants; their names in a design document are not evidence that every engine enforces them.

Several difficult cases remain visible rather than disappearing into the proof. An in-place tensor update needs an undo method or a preflight rule strong enough to prevent partial modification on the covered fault classes. A long matrix instruction that permits partial architectural completion needs a restart cursor or another precise description of which suboperations are committed. An error discovered after retirement, such as a later asynchronous integrity failure, falls outside the before-commit proof unless the architecture retains the necessary recovery information. Cancellation must establish whether a command has actually stopped, whether its responses can still arrive, and which resources it continues to own. Silently treating ``cancel requested'' as ``all effects canceled'' breaks both state recovery and allocation safety.

Remote publication is the hardest boundary. Suppose a younger instruction on A sends a value that B uses to update its own state, and an older instruction on A subsequently faults. Restoring A's rename map cannot undo B's computation. A design must therefore delay the send's visibility until the relevant local boundary is irrevocable, define the send as an explicit non-rollback boundary, or use a distributed protocol that can coordinate recovery. Delaying publication can retain a large output and add retirement-head blocking. A distributed protocol adds epochs, acknowledgment, duplicate suppression or idempotency, and agreement about what has committed. None follows from an instruction carrying a ROB identifier.

\subsection{NVIDIA GPU ownership and tooling}

PTX provides an explicit scoped memory model for the documented $\texttt{sm\_70}$-and-later scope, with specified exclusions, rather than treating asynchronous execution as unrestricted memory behavior. Its proxy rules also matter when asynchronous engines and ordinary accesses interact; ordinary ordering and asynchronous completion are separate obligations \citep{u007_ptx}. At the runtime layer, streams and events let software establish dependencies without synchronizing the entire host and device. A stream wait can defer downstream work while unrelated work remains eligible \citep{u007_cuda_async}.

NVIDIA's debugging facilities expose useful but different contracts. Compute Sanitizer distinguishes precisely attributed instrumented access errors from less precisely attributed hardware exceptions. Its race and synchronization checking have documented coverage; they do not prove arbitrary global or distributed code race-free, and error termination does not roll back earlier outputs \citep{u007_compute_sanitizer}. CUDA-GDB documents fault-origin reporting separately from the later stopping PC, and offers autostep to improve localization in selected regions \citep{u007_cuda_gdb}.

The algorithm or library chooses synchronization scope and buffer ownership; the compiler preserves the relevant source semantics; runtime events connect launches; device mechanisms implement the specified ordering and completion. The resulting control inventory includes per-thread or warp execution state, readiness and dependency tracking, barrier state, asynchronous command identities, and debug instrumentation where enabled. The exact circuit topology need not be known to recognize these functions. Conversely, profiler or debugger features do not establish a CPU-style ROB implementation or universal instruction restart.

The benefit condition is concrete: explicit scopes and event dependencies are valuable when independent work exists outside the required dependency cone. Broad waits may be simpler to audit but retain less overlap; finer dependencies require more precise ownership reasoning and state. Instrumentation can improve attribution while perturbing the schedule that exposed the original error. A fair comparison therefore records production semantics, diagnostic mode, and recovery scope separately rather than comparing an instrumented GPU with an uninstrumented proposal.

\subsection{Ascend compiler and framework synchronization}

Ascend C describes scalar issue to asynchronous vector, matrix, and movement pipelines. In the cited CANN 8.5 interface, cross-pipeline flags connect producer completion to consumer execution, while pipeline barriers constrain a selected pipeline. The compiler and framework insert substantial synchronization, including documented Vector and Cube responsibilities; manual cases remain when the automatic analysis or movement-address assumptions are insufficient. The documentation distinguishes portable TQueSync interfaces from hardware-sensitive ISASI forms \citep{u007_ascend_sync}. Consequently, describing Ascend kernels as universally hand-scheduled barriers would overstate the author burden, while treating framework synchronization as an unrestricted proof would hide its assumptions.

msSanitizer documents memory-hazard checking across supported memory and pipeline combinations, with source-context reporting and release-dependent coverage \citep{u007_ascend_sanitizer}. msDebug supports hardware debugging and core analysis, but explicitly warns that execution can continue for several instructions before an exception is reported: saved registers and memory can be inaccurate even when the PC is corrected \citep{u007_ascend_debug}. This is direct evidence that accurate attribution and coherent recoverable state must be assessed separately, without dismissing the diagnostic value of the tools.

Responsibility is distributed across algorithmic staging, tensor-lifetime analysis, framework-owned flags, runtime completion, and hardware pipelines. The design costs include event state and ownership, local buffer lifetimes, waits inserted conservatively when aliases are uncertain, and the resources used by diagnostic instrumentation. Automation is especially valuable when a kernel stays inside the framework's analyzable ownership model. Custom aliasing, unusual pipeline interactions, or library composition can shift obligations back to the author. The cited term ``committed'' in a pipeline-barrier description should be interpreted under that API, not promoted to an unprovided architectural retirement guarantee.

\subsection{TPU compiler dependencies and low-level Pallas control}

OpenXLA uses tokens to express side-effect dependencies, including joins through \texttt{AfterAll} \citep{u007_xla_tokens}. This is a compiler-IR contract. It does not reveal a complete physical TPU memory model or imply that independent physical engines execute sequentially. Pallas makes a lower-level boundary visible: software pipelining coordinates asynchronous units and buffer lifetimes, trading additional storage for earlier movement and later waits \citep{u007_pallas_pipeline}. The compiler can absorb regular scheduling work when the kernel supplies the necessary dependence and lifetime information.

For distributed kernels, the programmer still needs correctly matched transfers, synchronization counts, slot reuse, and participant identity. Pallas documents separate completion milestones and failure modes for mismatched synchronization \citep{u007_pallas_distributed}. Its debugging guide describes interpreted race detection and compiler-pass dumps \citep{u007_pallas_debug}. A simulation can explore memory and synchronization behavior and a compiler dump can expose a missing edge; neither is evidence of the exact silicon execution that failed, nor proof that all permitted schedules were checked.

The costs arise in both generated and explicit control: semaphore state, buffer versions, scheduling metadata, and the retained payload needed while communication overlaps computation. A compiler can simplify the normal programmer's contract without removing these functions. The proposed superscalar alternative should therefore be compared against the compiler's optimized schedule and actual debug facilities, not against an imagined TPU with no synchronization or visibility. Public evidence reviewed here leaves general instruction-level precise restart unspecified; that is an evidence boundary, not a claim that a particular hidden microarchitecture cannot support recovery.

\subsection{Tenstorrent explicit ownership and scoped communication}

TT-Metalium documents that writes on the same NoC and virtual channel follow program order; using different NoCs or virtual channels does not automatically preserve that order \citep{u007_tt_ordering}. Reader, compute, and writer roles must connect movement completion to circular-buffer publication and consumption. A queue-ready condition says that the producer has published data under the queue contract, so the producer must establish the required payload completion first. The receiver's eventual slot release is another event.

Watcher provides kernel identities, per-controller waypoints, and selected address or bounds checks. Its current hardware-exception PC/cause reporting is explicitly Quasar-specific \citep{u007_tt_watcher}; that capability cannot be silently attributed to Wormhole or Blackhole. The experimental NoC Debug Dump reports selected missing barriers and outstanding writes, while warning that timing can hide faults and tracing adds overhead \citep{u007_tt_noc_debug}. Waypoints and traces are valuable localization aids, but do not constitute a coherent global checkpoint.

The ownership split is particularly visible: kernels and libraries arrange roles, channels, buffer handoff, and completion; controllers and engines advance the specified operations; runtime or fabric software manages broader program lifetime. The benefit is that communication and compute control can be separated and overlapped with explicit storage boundaries. The price includes controller instructions, buffer and semaphore state, trace capacity, and the need to compose several local protocols correctly. This is a different distribution of control from a superscalar instruction window, not an absence of control logic.

\subsection{The proposed superscalar NPU and the full cost of precision}

The supplied project inspection records ROB control, renamed identities, completion and non-flush gates in a timing model, with a separate typed bring-up path containing its own identity and recovery contracts. These observations motivate a local precise-state design; they do not establish one integrated, verified implementation. In particular, neither a model component nor a source-level retirement routine proves that every tensor update, local transfer, atomic, and remote effect obeys the invariants above.

The proposal's strongest plausible advantage is a simpler local programming and debugging boundary: source-level operations could retain stable provenance while dependency tracking permits safe internal overlap. That advantage applies where the architecture can identify dependencies and control publication, and where the cost of retaining recoverable state is acceptable. It becomes weaker for irrevocable remote effects, uncertain memory aliases, very large in-place updates, or long-latency operations that block the retirement frontier. These cases need explicit semantics, rather than a blanket claim that out-of-order hardware makes asynchronous code easy to debug.

A useful provenance chain maps source operation to compiler IR, emitted instruction or command, dynamic epoch, engine issue, completion, retirement, and fault. Many-to-one fusion and one-to-many lowering must remain representable. Trace records should identify lost records and distinguish per-engine order from cross-engine timestamp order. Otherwise a visually tidy timeline can suggest causal edges that the machine never guaranteed. Deterministic replay, if desired, needs a separate account of external inputs and scheduling decisions that influence execution.

The cost ledger must include exception records, tags and generations, committed and speculative mappings, retained payload versions, output-publication buffers when needed, and arbitration on recovery paths. Retaining an old tile can occupy an existing physical pool rather than require a new full-tile copy; the important quantity is unavailable capacity over time, not only copy count. Let $b_v$ be the bytes occupied by version $v$ and $[a_v,r_v)$ its retention interval. The analytical occupancy integral
\begin{equation}
 \mathcal{R}=\sum_v b_v(r_v-a_v)
\end{equation}
helps compare recovery schemes with different lifetimes, while peak occupancy determines capacity feasibility. Neither quantity alone predicts energy, area, or throughput. Metadata width, array ports, selection logic, checkpoint traffic, and the critical timing path require implementation evidence.

The appropriate correctness result is therefore a stated recovery domain and a proof or verification argument for every effect that can cross it. The appropriate performance result is reduced unnecessary ordering or debugging work after counting the storage and control that preserve that domain. Local retirement, scoped memory order, asynchronous completion, and distributed recovery remain separate contracts whose composition must be designed deliberately.
